\renewcommand\thechapter{c2.9a}
\lab{Electrical measurements}\label{lab:covid-electrical-measurements}

\section*{About this lab}

\covid\ 
It is intended to be used during the ninth week in Physics 222.

\apparatus
\equip{1.5 V battery}
\equip{battery holder}
\equip{banana plug cables (6)}
\equip{alligator clips (4)}
\equip{flashlight bulb}
\equip{bulb holder}
\equip{multimeter}

\begin{goals}

\item[] Build a complete circuit and light a lightbulb.

\item[] Create and observe an open circuit and a short circuit.

\item[] Measure a voltage and a current.

\item[] Determine the resistance of a lightbulb.

\end{goals}

\labpart{Lighting a lightbulb}

The figure below shows a complete circuit in which energy supplied
by a battery is used to light up a lightbulb. 

\fig{em-mea-covid-basic-circuit}

The second figure shows a more abstract representation of this circuit,
called a schematic. It shows what's connected to what.

\fig{em-mea-covid-basic-circuit-schematic}

Build the circuit as shown. The setup in the photo was built using the
materials in your kit. The connections to the lightbulb holder are most
easily made by clipping alligator clips sideways onto the screws.

Throughout this lab, you do not have to worry about shocking yourself.
The 1.5 V battery doesn't have a high enough voltage to hurt you.
(Even a 9 V battery won't do anything to you when you handle it,
although an important educational experience for little kids is
to touch one to their tongue, which causes a tingling sensation.)

\labpart{Open circuit, short circuit}

Observe what happens when you make an open circuit or a short circuit, as shown
in the schematics. When you do the short circuit, don't leave the extra wire in
place for too long, or you will create a lot of heat and possibly burn your fingers.

\fig{em-mea-covid-basic-circuit-open-and-short}

\labpart{Measuring voltage and current}


As shown in the figure, measuring current and voltage (potential difference) requires
hooking the meter into the circuit in two completely different
ways. This highlights the fact that they are completely different
physical quantities. Current measures the flow of electric charge,
while voltage measures a potential energy difference per unit charhe.
You have already had the experience of measuring a voltage difference
(difference in electrical potential) in lab \ref{lab:electric-field}.

\fig{em-mea-covid-i-and-v-measurements}

The arrangement for the ammeter is called a series circuit,
because every charged particle that travels the circuit has
to go through each component, one after another.
The series circuit is arranged like beads on a necklace.

The setup for the voltmeter is an example of a parallel
circuit. A charged particle flowing, say, clockwise
around the circuit passes through the power supply and
then reaches a fork in the road, where it has a choice
of which way to go. Some particles will pass through the
bulb, others (not as many) through the meter; all of them
are reunited when they reach the junction on the right.

Students tend to have a mental block against setting up the
ammeter correctly in series, because it involves breaking the circuit
apart in order to insert the meter.  If you hook up the ammeter incorrectly,
in parallel rather than in series, the meter provides an easy path for the flow
of current, so a large amount of current will flow. To protect the meter
from this surge, there is a fuse inside, which will blow, and producing
an open circuit. This is not a huge tragedy, but you will have to recognize
when it happens and know to replace the fuse. Replacement fuses are included
in your kit.

We have a single meter that plays both the role of the voltmeter
and the role of the ammeter in this lab. Because it can do both
these things, it is referred to as a multimeter. Multimeters are
highly standardized, and the following instructions are generic
ones that will work with whatever meters you happen to be using
in this lab.

\emph{Measuring a voltage}

Two wires connect the meter to the circuit. 
At the places where three wires come together at
one point, you can plug a banana plug into the back of
another banana plug. At the meter, make one connection
at the ``common'' socket (``COM'') and the other at the
socket labeled ``V'' for volts. The common plug is called
that because it is used for every measurement, not just for
voltage.

Many meters, including the one we're planning on sending out
with our kits, have both DC and AC voltage settings. The battery is
a DC voltage source, meaning that its voltage has a fixed polarity rather
than alternating some number of times per second. The DC voltage
settings are marked with this hieroglyph \raisebox{-.2\height}{\includegraphics{../figs/hieroglyph-dc}},
AC with this one \raisebox{-.2\height}{\includegraphics{../figs/hieroglyph-ac}}.

Many multimeters have more than one scale for
measuring a given thing. For instance, a meter may have a
millivolt scale and a volt scale. One is used for measuring
small voltage differences and the other for large ones. You may not
be sure in advance what scale is appropriate, but that's not
a big problem --- once everything is hooked up, you can try
different scales and see what's appropriate. Use the switch
or buttons on the front to select one of the voltage scales.
By trial and error, find the most precise scale that doesn't
cause the meter to display an error message about being overloaded.

Write down your measurement, with the units of volts, and stop
for a moment to think about what it is that you've measured.
Imagine holding your breath and trying to make your eyeballs
pop out with the pressure. 
Intuitively, the voltage difference is like the pressure difference
between the inside and outside of your body.

What do you think will happen if you unscrew the bulb, leaving
an air gap, while the
power supply and the voltmeter are still going? Try it.
Interpret your observation in terms of the breath-holding
metaphor.

\emph{Current}

The procedure for measuring the current differs only because
you have to hook the meter up in series and because you have
to use the ``A'' (amps) plug on the meter and select a current
scale.

In the breath-holding metaphor, the number you're measuring
now is like the rate at which air flows 
through your lips as you let it hiss out. Based on this
metaphor, what do you think will happen to the reading when
you unscrew the bulb? Try it.

\labpart{Resistance}

The ratio of voltage difference to current is called the resistance
of the bulb, $R=\Delta V/I$. Its units of volts per amp can be
abbreviated as ohms, $\Omega$ (capital Greek letter omega).

Calculate the resistance of your lightbulb. Resistance is the
electrical equivalent of kinetic friction. Just as rubbing
your hands together heats them up, objects that have
electrical resistance produce heat when a current is passed
through them. This is why the bulb's filament gets hot enough
to heat up.

When you unscrew the bulb, leaving an air gap, what is the
resistance of the air?


